\subsection{Evaluation Methodology}\label{sec:evaluation}
$G$ was evaluated on the test and validation splits. Each row of the window matrix consists of a condition window of $l$ realized returns and the realized return of the next day. The conditions were passed to the trained $G$ and the generated values were compared to the realized values. For each condition, $B=1000$ samples were drawn from $G$ with different noise vectors. 

The realized returns of the split, $x = (x_1, \ldots, x_n)$, give the reference values, where $n$ is the number of rows in the window matrix. 
The generated draws form a matrix $\hat{X} \in \mathbb{R}^{n \times B}$, where $\hat{x}_{i,j}$ is the $j$-th draw of condition $i$. Additional sets were derived from $\hat{X}$. 

The path set $\hat{x}_{\mathrm{path}} = (\hat{x}_{1,1}, \ldots, \hat{x}_{n,1})$ is the first column of $\hat{X}$, giving one generated return per trading day.
The pooled set takes all $nB$ generated entries of $\hat{X}$.

\begin{equation}
    \hat{x}_{\mathrm{pool}} = \{\, \hat{x}_{i,j} \;|\;
    i = 1, \ldots, n, \; j = 1, \ldots, B \,\}.
\end{equation}
 
The mean set holds the mean of all draws for each condition

\begin{equation}
    m = (m_1, \ldots, m_n), 
    \text{ where }
    m_i = \frac{1}{B} \sum_{j=1}^{B} \hat{x}_{i,j}.
\end{equation}


$\hat{x}_{\mathrm{path}}$ is directly comparable with $x$, since both sets have the same length. $\hat{x}_{\mathrm{pool}}$ contains $B$ times more observations than $x$. 

Distributional statistics were computed for $x$, $\hat{x}_{\mathrm{path}}$ and $\hat{x}_{\mathrm{pool}}$: mean, standard deviation, skewness, excess kurtosis, minimum, maximum, $1\%$ and $99\%$ quantiles. Additionally, the fractions of returns exceeding $5\%$ and
$10\%$ in absolute value were computed.

The point forecast accuracy was measured using $x$ and $m$. The mean absolute error (MAE), the root mean squared error (RMSE) and correlation were computed, following \textcite{vuleticFinGANForecastingClassifying2024}. The fraction of positive values in $m$ was calculated to check whether the
forecasts are biased in one direction. The standard deviation of $m$ was calculated to check how much the forecasts vary across conditions. 

The autocorrelation of the squared returns was computed for $x$ and
$\hat{x}_{\mathrm{path}}$ at lags $1$, $2$, $5$ and $10$. The autocorrelation measures whether large returns are followed by
large returns. $\hat{x}_{\mathrm{pool}}$ cannot be used since it has no order in time.

The weighted strategy described in Sec.~\ref{sec:fin-gan} was applied to the rows of $\hat{X}$ and $x$. $\overline{\mathrm{wPnL}}$ and $\mathrm{SR}_{w}$ were computed. 

Lastly, the mode collapse described in Sec.~\ref{sec:fin-gan} was checked. For distribution mode collapse, the standard deviation of each row of $\hat{X}$ was computed. The smallest value was compared with $\epsilon$. For mean mode collapse, the standard deviation of $m$ was compared with $\epsilon$.

% TODO: add motivation for the evaluation metrics once Sec. 2.3 exists.
% - distributional statistics: which stylized facts they measure
% - exceedance fractions: link to OTM option value
% - ACF of squared returns: why squared, volatility clustering as a
%   stylized fact
% - MAE/RMSE/correlation: comparability with the equity results of
%   Vuletic et al., not the target of the model
% - SR_w: directional accuracy only, cross-ref 2.5.4